\documentclass[a4paper,14pt]{extarticle}

\usepackage{polyglossia}
\setdefaultlanguage{russian}

\usepackage{amsmath, amssymb}
\usepackage{graphicx}
\usepackage{geometry}
\usepackage{setspace}
\usepackage{indentfirst}
\usepackage{fontspec}
\usepackage{tabularx}
\usepackage{colortbl}
\usepackage{pgfplots}

\pgfplotsset{compat=1.18}
\definecolor{lightgray}{gray}{0.9}
\setmainfont{Times New Roman}

\geometry{left=2cm,right=2cm,top=2cm,bottom=2cm}
\onehalfspacing
\setlength{\parindent}{1.25cm}

\begin{document}

\begin{titlepage}
\begin{center}
\textbf{Министерство образования Республики Беларусь}\\
\textbf{Белорусский государственный университет}

Факультет радиофизики и компьютерных технологий
\vspace{6cm}

\textbf{Лабораторная работа №5}\\
\textbf{«Зависимость электропроводности от температуры»}
\end{center}

\vfill
\begin{flushright}
Выполнила:\\
Студентка 1 курса 7 группы\\
Данильченко Мария\\
Преподаватель: Садов П. В.\\
\end{flushright}

\vfill
\begin{center}
Минск, 2026
\end{center}
\end{titlepage}

\section*{Цель работы}

Изучить характер зависимости сопротивления металлов и полупроводников от температуры. Определить температурный коэффициент сопротивления металла и характер температурной зависимости проводимости полупроводника.

% -------------------- Оборудование --------------------
\section*{Оборудование и приборы}

Сушильный шкаф, термометр и датчик температуры, измеритель сопротивления.

% -------------------- Теория --------------------
\section*{Теоретическая часть}

Удельное сопротивление различных металлов при комнатной температуре имеет значения
$10^{-8} \div 10^{-6}\,\Omega\cdot\text{м}$.
Твёрдые вещества с удельным сопротивлением
$(10^{10} \div 10^{20}\,\Omega\cdot\text{м})$
являются диэлектриками. Вещества с промежуточными значениями
$(10^{-4} \div 10^{10}\,\Omega\cdot\text{м})$
называются полупроводниками.

С повышением температуры сопротивление металлов увеличивается, тогда как у полупроводников уменьшается. Это связано с различными механизмами переноса заряда.

Для чистых металлов при $T > 273\,\text{К}$ зависимость сопротивления от температуры описывается выражением
\[
R = R_0(1 + \alpha t + \beta t^2).
\]

При небольших изменениях температуры используется линейное приближение, позволяющее определить температурный коэффициент сопротивления.

Для полупроводников сопротивление определяется концентрацией носителей заряда, которая экспоненциально зависит от температуры:
\[
R = R_0 \exp\!\left(\frac{\Delta E}{2kT}\right),
\]
где $\Delta E$ — ширина запрещённой зоны, $k = 1{,}38 \times 10^{-23}$\,Дж/К — постоянная Больцмана, $T$ — абсолютная температура. Логарифмируя, получаем линейную зависимость:
\[
\ln R = \frac{\Delta E}{2k} \cdot \frac{1}{T} + \ln R_0.
\]
Таким образом, на графике $\ln R$ от $1/T$ экспериментальные точки должны лечь на прямую с наклоном
\[
\operatorname{tg}\varphi = \frac{\Delta E}{2k},
\]
откуда ширина запрещённой зоны:
\[
\Delta E = 2k \cdot \operatorname{tg}\varphi.
\]

% -------------------- Практика --------------------
\section*{Практическая часть}

\subsection*{Экспериментальная установка}

Исследуемые образцы помещались в сушильный шкаф. Температура контролировалась термометром, сопротивление измерялось омметром.

\subsection*{Результаты измерений}

\begin{table}[h]
\centering
\caption{Результаты измерений}
\begin{tabular}{|c|c|c|c|}
\hline
$t,^\circ$C & $R_1$, Ом & $R_2$, Ом & $R_3$, Ом \\
\hline
33  & 1.142 & 1.192 & 1.480 \\
43  & 1.178 & 1.306 & 0.950 \\
53  & 1.209 & 1.403 & 0.750 \\
63  & 1.240 & 1.500 & 0.530 \\
73  & 1.271 & 1.603 & 0.400 \\
83  & 1.300 & 1.707 & 0.333 \\
93  & 1.325 & 1.718 & 0.312 \\
103 & 1.343 & 1.743 & 0.300 \\
\hline
\end{tabular}
\end{table}

\subsection*{Графики зависимости сопротивления от температуры}

\begin{figure}[h]
\centering
\begin{tikzpicture}
\begin{axis}[
    width=0.9\textwidth,
    xlabel={$t,^\circ$C},
    ylabel={$R_1$, Ом},
    grid=both,
    title={Первый образец}
]
\addplot[only marks, mark=*] table {
33 1.142
43 1.178
53 1.209
63 1.240
73 1.271
83 1.300
93 1.325
103 1.343
};
\addplot[domain=30:105, samples=100]{0.00291*x + 1.053};
\end{axis}
\end{tikzpicture}
\caption{Зависимость сопротивления первого образца от температуры}
\end{figure}

\begin{figure}[h]
\centering
\begin{tikzpicture}
\begin{axis}[
    width=0.9\textwidth,
    xlabel={$t,^\circ$C},
    ylabel={$R_2$, Ом},
    grid=both,
    title={Второй образец}
]
\addplot[only marks, mark=*] table {
33 1.192
43 1.306
53 1.403
63 1.500
73 1.603
83 1.707
93 1.718
103 1.743
};
\addplot[domain=30:105, samples=100]{0.00825*x + 0.960};
\end{axis}
\end{tikzpicture}
\caption{Зависимость сопротивления второго образца от температуры}
\end{figure}

\begin{figure}[h]
\centering
\begin{tikzpicture}
\begin{axis}[
    width=0.9\textwidth,
    xlabel={$t,^\circ$C},
    ylabel={$R_3$, Ом},
    grid=both,
    title={Третий образец (полупроводник)}
]
\addplot[only marks, mark=*] table {
33 1.480
43 0.950
53 0.750
63 0.530
73 0.400
83 0.333
93 0.312
103 0.300
};
\addplot[domain=30:105, samples=100]{-0.0153*x + 1.671};
\end{axis}
\end{tikzpicture}
\caption{Зависимость сопротивления полупроводника от температуры}
\end{figure}

% -------------------- Обработка результатов --------------------
\section*{Обработка результатов. Определение $\Delta E$ полупроводника}

Для определения ширины запрещённой зоны полупроводника (третий образец) построим
график $\ln R_3$ от $1/T$. Значения $T = t + 273{,}15\,\text{К}$, вычисленные
величины $1/T$ и $\ln R_3$ приведены в таблице~2.

\begin{table}[h]
\centering
\caption{Данные для графика $\ln R_3$ от $1/T$}
\begin{tabular}{|c|c|c|c|c|}
\hline
$t$, °C & $T$, К & $1/T \times 10^{3}$, К$^{-1}$ & $R_3$, Ом & $\ln R_3$ \\
\hline
33  & 306,1 & 3,266 & 1,480 &  0,392 \\
43  & 316,1 & 3,163 & 0,950 & $-$0,051 \\
53  & 326,1 & 3,066 & 0,750 & $-$0,288 \\
63  & 336,1 & 2,975 & 0,530 & $-$0,635 \\
73  & 346,1 & 2,889 & 0,400 & $-$0,916 \\
83  & 356,1 & 2,808 & 0,333 & $-$1,100 \\
93  & 366,1 & 2,731 & 0,312 & $-$1,165 \\
103 & 376,1 & 2,659 & 0,300 & $-$1,204 \\
\hline
\end{tabular}
\end{table}

\begin{figure}[h]
\centering
\begin{tikzpicture}
\begin{axis}[
    width=0.9\textwidth,
    xlabel={$1/T \times 10^{3}$, К$^{-1}$},
    ylabel={$\ln R_3$},
    grid=both,
    title={График $\ln R_3$ от $1/T$ для полупроводника},
    xtick distance=0.1,
    xticklabel style={/pgf/number format/fixed, /pgf/number format/precision=3}
]
\addplot[only marks, mark=*] coordinates {
(3.266,  0.392)
(3.163, -0.051)
(3.066, -0.288)
(2.975, -0.635)
(2.889, -0.916)
(2.808, -1.100)
(2.731, -1.165)
(2.659, -1.204)
};
\addplot[domain=2.6:3.3, samples=100]{2.7051*x - 8.5863};
\end{axis}
\end{tikzpicture}
\caption{Линеаризованная зависимость $\ln R_3$ от $1/T$}
\end{figure}

\subsection*{Вычисление $\operatorname{tg}\varphi$ и $\Delta E$}

Методом наименьших квадратов по данным таблицы~2 получена прямая:
\[
\ln R_3 = \frac{\Delta E}{2k} \cdot \frac{1}{T} + b,
\]
коэффициент наклона которой:
\[
\operatorname{tg}\varphi = \frac{\Delta E}{2k} = 2705{,}1\,\text{К}.
\]

Ширина запрещённой зоны:
\[
\Delta E = 2k \cdot \operatorname{tg}\varphi = 2 \times 1{,}38 \times 10^{-23}\,\frac{\text{Дж}}{\text{К}}
           \times 2705{,}1\,\text{К}
         = 7{,}47 \times 10^{-20}\,\text{Дж}.
\]

Переводим в электрон-вольты ($1\,\text{эВ} = 1{,}6 \times 10^{-19}\,\text{Дж}$):
\[
\boxed{\Delta E \approx 0{,}47\,\text{эВ}.}
\]

% -------------------- Вывод --------------------
\section*{Вывод}

В ходе лабораторной работы была исследована температурная зависимость сопротивления
трёх образцов. Установлено, что первый и второй образцы являются металлами, так как
их сопротивление увеличивается с ростом температуры. Третий образец проявляет
свойства полупроводника, поскольку его сопротивление уменьшается при нагревании.

Для полупроводника по линеаризованному графику $\ln R_3$ от $1/T$ определены
тангенс угла наклона $\operatorname{tg}\varphi = 2705{,}1\,\text{К}$ и ширина
запрещённой зоны $\Delta E \approx 0{,}47\,\text{эВ}$.
Полученные экспериментальные результаты согласуются с теоретическими представлениями
о природе проводимости полупроводников.

\end{document}
